% ============================================================
%  preamble.tex — 日本語数学ノート テンプレート共通設定
%
%  定理環境のカウンタ規約は tools/site が再現している：
%    - 全定理環境が definition のカウンタを共有（use counter from）
%    - カウンタは節ごとにリセット（number within=section）
%  この規約を変えるときは tools/site/lib/tex2md.mjs の
%  buildNumberMap() も合わせること（PDF とサイトで番号がずれる）。
% ============================================================

% --- 基本パッケージ ---
\usepackage{amsmath,amssymb,amsthm,mathtools}
\usepackage[dvipdfmx,hidelinks]{hyperref}
\usepackage{pxjahyper}

% --- 図・表 ---
\usepackage{tikz}
\usepackage{float}
\usetikzlibrary{decorations.pathreplacing,calc,arrows.meta,patterns,positioning}
\usepackage{pgfplots}
\pgfplotsset{compat=1.18}
\usepackage{graphicx}
\graphicspath{{fig/}}

% --- 定理環境 (tcolorbox) ---
% 第 2 引数はラベル接頭辞。tools/site はこの接頭辞で \ref を解決するので、
% def / prop / thm / clm / lem / cor / rem / ex は変えないこと。
\usepackage{tcolorbox}
\tcbuselibrary{theorems,breakable}

% Definition: 青系
\newtcbtheorem[number within=section]{problem}{問題}{
  colback=teal!4, colframe=teal!55!black, fonttitle=\bfseries, breakable
}{prob}
% Definition: 青系
\newtcbtheorem[number within=section]{definition}{Def.}{
  colback=blue!5, colframe=blue!40!black, fonttitle=\bfseries, breakable
}{def}
% Proposition: オレンジ系
\newtcbtheorem[use counter from=definition]{proposition}{Prop.}{
  colback=orange!5, colframe=orange!50!black, fonttitle=\bfseries, breakable
}{prop}
% Theorem: 赤系（章・節の目玉となる重要結果）
\newtcbtheorem[use counter from=definition]{theorem}{Thm.}{
  colback=red!5, colframe=red!40!black, fonttitle=\bfseries, breakable
}{thm}
% Claim: 緑系（Thm ほど深くはないが，証明を要する主張）
\newtcbtheorem[use counter from=definition]{claim}{Claim}{
  colback=green!5, colframe=green!45!black, fonttitle=\bfseries, breakable
}{clm}
% Lemma: シアン系（命題の証明に用いる補助結果）
\newtcbtheorem[use counter from=definition]{lemma}{Lem.}{
  colback=cyan!5, colframe=cyan!50!black, fonttitle=\bfseries, breakable
}{lem}
% Corollary: すみれ系（定理から直ちに従う結果）
\newtcbtheorem[use counter from=definition]{corollary}{Cor.}{
  colback=violet!5, colframe=violet!40!black, fonttitle=\bfseries, breakable
}{cor}
% Remark: 黄系
\newtcbtheorem[use counter from=definition]{remark}{Rem.}{
  colback=yellow!5, colframe=yellow!40!black, fonttitle=\bfseries, breakable
}{rem}
% Example: 灰色系
\newtcbtheorem[use counter from=definition]{example}{Ex.}{
  colback=gray!5, colframe=gray!50!black, fonttitle=\bfseries, breakable
}{ex}
% Memo: 薄い黒系（番号なし，発展的補足）
\newtcolorbox{memo*}{
  colback=black!7, colframe=black!40,
  fonttitle=\bfseries, title={Memo},
  breakable
}
% 証明の直後が enumerate でも、最初の項目だけ見出しの横へずれないようにする。
% 通常の文章で始まる証明は従来どおり見出しと同じ行から始まる。
\makeatletter
\renewenvironment{proof}[1][\proofname]{\par
  \pushQED{\qed}%
  \normalfont \topsep6\p@\@plus6\p@\relax
  \trivlist
  \item[\hskip\labelsep\itshape #1\@addpunct{.}]\ignorespaces\leavevmode
}{%
  \popQED\endtrivlist\@endpefalse
}
\makeatother

% 証明環境の装飾
\tcolorboxenvironment{proof}{
  colback=white, colframe=black!20,
  leftrule=2pt, rightrule=0pt, toprule=0pt, bottomrule=0pt,
  breakable
}

% --- アルゴリズム環境 ---
\newcounter{algcounter}[section]
\newtcolorbox{algorithm}[1]{
  colback=white, colframe=black!70,
  fonttitle=\bfseries,
  title={Algorithm~\refstepcounter{algcounter}\thealgcounter\label{#1}},
  boxrule=0.8pt
}

% --- サイト向けの指示マクロ（PDF では何も出力しない）---
% \demohint{名前}: site.config.mjs の demos に登録した図をサイト側で差し込む
\newcommand{\demohint}[1]{}
% \blockmeta{...}: ブロックの機械可読メタデータ。PDF にもサイトにも出さない
\newcommand{\blockmeta}[1]{}

% ============================================================
%  数式マクロ
%  ここに書いたものは tools/site が自動抽出して MathJax にも渡す。
%  サイト側の設定に書き写す必要はない。
%  MathJax が解釈できない綴りだけ site.config.mjs の macroOverrides で
%  差し替えること（例: stmaryrd の \llbracket）。
% ============================================================

\newcommand{\R}{\mathbb{R}}
\newcommand{\N}{\mathbb{N}}
\newcommand{\Z}{\mathbb{Z}}
\newcommand{\Q}{\mathbb{Q}}
\newcommand{\E}{\mathbb{E}}

\newcommand{\X}{\mathcal{X}}
\newcommand{\Y}{\mathcal{Y}}
\newcommand{\Bb}{\mathcal{B}}
\newcommand{\Prob}{\mathcal{P}}

\DeclareMathOperator*{\argmin}{arg\,min}
\DeclareMathOperator*{\argmax}{arg\,max}
\DeclareMathOperator{\diag}{diag}
\DeclareMathOperator{\tr}{tr}
\DeclareMathOperator{\supp}{supp}
\DeclareMathOperator{\Id}{Id}

\newcommand{\abs}[1]{\lvert #1 \rvert}
\newcommand{\norm}[1]{\lVert #1 \rVert}
\newcommand{\inner}[2]{\langle #1,\, #2 \rangle}
\newcommand{\defeq}{\overset{\mathrm{def}}{=}}
\renewcommand{\d}{\mathrm{d}}

% --- 参考文献 ---
\usepackage[numbers,sort&compress]{natbib}
